% ---------------------------------------------------------------------------
% SECTION 7: TEMPORAL TRANSFER (NHANES 2015-2016)  [new second application]
% ---------------------------------------------------------------------------
\section{Temporal Transfer of the Calibrated Error Variance: NHANES
2015--2016}\label{sec:transfer}

A practitioner applying the method in real time cannot calibrate the
error variance on the data at hand if the validation measurement is
only collected later. This section therefore asks the question that
matters for practice: does the error variance calibrated on one survey
cycle transfer to another? We re-estimate the application of
Section~\ref{sec:app} on the previous NHANES cycle, 2015--2016
(DEMO\_H, BMX\_H, WHQ\_H, RXQ\_RX\_H), with the identical outcome
(number of prescription medications), the identical error-prone
covariate (self-reported weight), and the identical specification, but
we feed the methods the error variance calibrated on the 2017--2018
cycle ($\hat\sigma^2_\epsilon = 24.1\,\mathrm{kg}^2$) instead of
estimating it within the 2015--2016 sample. The within-cycle
calibration is available here as a benchmark: the 2015--2016
self-report minus measured differences have mean $-0.69$\,kg and
variance $21.8\,\mathrm{kg}^2$, so the transferred variance is about
$10\%$ larger than the cycle would calibrate on its own.

Table~\ref{tab:transfer} and Figure~\ref{fig:transfer} report the
results for $n = 6{,}100$. Three facts stand out. First, the ranking
and spacing of methods reproduces the 2017--2018 application: the
naive estimate is $+1.445$ per 10\,kg, the corrections cluster between
$+1.474$ (argmin) and $+1.526$ (regression calibration, corrected score,
standard SIMEX, and A-SIMEX 1-SE), and the measured-weight anchor is $+1.323$, below the
corrected estimates as in Section~\ref{sec:app}, consistently with the
nonclassical component of self-report error. Second, the transferred
and within-cycle calibrations deliver essentially identical estimates:
the largest gap across methods and variance inputs is
$0.009$, well under one standard error ($0.083$). A $10\%$
misstatement of the error variance moves the corrected estimates by a
negligible amount, exactly as the sensitivity analysis of
Section~\ref{sec:app-sens} predicted. Third, the CV layer again
selects minimal extrapolation and remains stable across cycles.

The practical reading for applied work: a laboratory that has
calibrated its measurement error once (here, a national survey cycle)
can carry that calibration to the next wave, and the correction
inherits the calibration error only in a attenuated form. The
temporal stability of $\hat\sigma^2_\epsilon$ itself ($24.1$ on the
2017--2018 complete-case sample versus $21.8$ here) is a property of
the survey instrument rather than of the method, and its drift is
visible in the small difference between the transferred and
within-cycle columns.

\begin{table}[t]
\centering
\small
\caption{Temporal transfer, NHANES 2015--2016 ($n=6100$):
Poisson regression of the number of prescription medications on
self-reported weight (coefficient per 10\,kg, robust standard
errors). The transferred column feeds the methods the error
variance calibrated on the 2017--2018 cycle
($\hat\sigma^2_\epsilon = 24.1\,\mathrm{kg}^2$);
the within column re-estimates it from the 2015--2016
self-report minus measured differences
($\hat\sigma^2_\epsilon = 21.8\,\mathrm{kg}^2$).}
\label{tab:transfer}
\begin{tabular}{lcccc}
\toprule
 & \multicolumn{2}{c}{Transferred $\hat\sigma^2_\epsilon$} & \multicolumn{2}{c}{Within-cycle $\hat\sigma^2_\epsilon$} \\
\cmidrule(lr){2-3}\cmidrule(lr){4-5}
Method & $\hat\beta$ & SE & $\hat\beta$ & SE \\
\midrule
Naive (self-reported weight) & +1.445 & (0.079) & +1.445 & (0.079) \\
Regression calibration & +1.526 & (0.084) & +1.518 & (0.083) \\
Corrected score (Nakamura) & +1.526 & (0.083) & +1.518 & (0.082) \\
Standard SIMEX & +1.523 & (0.083) & +1.518 & (0.083) \\
Adaptive SIMEX (1-SE rule) & +1.523 & (0.083) & +1.514 & (0.082) \\
Adaptive SIMEX (argmin) & +1.474 & (0.081) & +1.475 & (0.081) \\
Measured weight (anchor) & +1.323 & (0.077) & +1.323 & (0.077) \\
\bottomrule
\end{tabular}
\end{table}

\begin{figure}[t]
\centering
\includegraphics[width=.8\textwidth]{fig8_transfer.pdf}
\caption{NHANES 2015--2016: weight coefficients (per 10\,kg, 95\% CIs)
with the transferred (green circles) and within-cycle (purple squares)
error variance. Dashed: measured weight (anchor); dotted: naive.}
\label{fig:transfer}
\end{figure}
